\documentclass[conference]{IEEEtran}

\usepackage{cite}
\usepackage{amsmath,amssymb,amsfonts}
\usepackage{algorithmic}
\usepackage{graphicx}
\usepackage{textcomp}
\usepackage{xcolor}
\usepackage{booktabs}
\usepackage{tagging}

\usepackage{sc26repro}

% For the FINAL submission the explanation/example tags are disabled so only the
% author-provided content renders. Uncomment to see the template guidance again.
% \usetag{explanation}
% \usetag{example}

\begin{document}

\twocolumn[%
{\begin{center}
\Huge
Appendix: Artifact Description/Artifact Evaluation
\end{center}}
]

%%%%%%%%%%%%%%%%%%%%%%%%%%%%%%%%%%%%%%%%%%%%%%%%%%%%%
%  AD Appendix
%%%%%%%%%%%%%%%%%%%%%%%%%%%%%%%%%%%%%%%%%%%%%%%%%%%%%

\appendixAD

\section{Overview of Contributions and Artifacts}

\subsection{Paper's Main Contributions}

\begin{description}
\item[$C_1$] \textbf{GGap}, a GPU-accelerated, agent-based re-implementation of the
  UVAFME forest gap model that preserves full ecological fidelity. A Site--Gap--Tree
  agent hierarchy with an 11-priority kernel pipeline resolves the model's sequential
  intra-year dependencies (shared nitrogen pool, 365-day soil biogeochemistry, light
  competition) on the GPU via CuPy JIT kernels on the SAGESim framework.
\item[$C_2$] A data-movement-optimized multi-GPU distribution: a Site--Gap--Tree
  \emph{colocation} invariant makes inter-rank MPI traffic scale with the number of
  \emph{sites} rather than trees, enabling GPU-aware MPI seed dispersal between
  geographically separated sites and strong/weak scaling on OLCF Frontier.
\item[$C_3$] Deterministic reproducibility (Philox RNG, partition-invariant) validated
  bitwise against the GAPpy CPU reference, together with a CONUS-scale ecological
  application producing species-composition and tree-size-distribution results across
  representative sites.
\end{description}

\subsection{Computational Artifacts}

\begin{description}
\item[$A_1$] GGap TN-region reproduction package. DOI: \texttt{https://doi.org/TODO}
  \emph{(TODO: mint Zenodo DOI for the archived snapshot of this repository).}
\end{description}

Artifact $A_1$ is a self-contained, downscaled reproduction (directory
\texttt{SC2026/tn\_example/}) that runs the GGap model on a two-GPU commodity node and
regenerates the science figure for one of the paper's ten representative sites. It
exercises $C_1$ (the full GPU model), $C_2$ (multi-GPU distribution with cross-rank
dispersal), and $C_3$ (deterministic output and ecological validation).

\begin{center}
\begin{tabular}{rll}
\toprule
Artifact ID  &  Contributions &  Related \\
             &  Supported     &  Paper Elements \\
\midrule
$A_1$   &  $C_1, C_2, C_3$ & Fig.~(10 representative \\
        &                  & sites), panel \#7 \\
        &                  & (S.~Appalachia) \\
\bottomrule
\end{tabular}
\end{center}


%%%%%%%%%%%%%%%%%%%%%%%%%%%%%%%%%%%%%%%%%%%%%%%%%%%%%%%%
\section{Artifact Identification}
%%%%%%%%%%%%%%%%%%%%%%%%%%%%%%%%%%%%%%%%%%%%%%%%%%%%%%%%

\newartifact

\artrel

Artifact $A_1$ is the reproduction package for contributions $C_1$--$C_3$. Because a
full-CONUS Frontier run is impractical for a reviewer, $A_1$ reduces the \emph{scale}
(a 54-site Tennessee sub-region) while keeping full per-site fidelity (500 gaps
$\times$ 1000 tree slots, 1000 simulation years). It runs the identical model and
distribution code path used for the paper's production runs, on two commodity GPUs.

\artexp

Running $A_1$ produces, for every site, a per-species time series of biomass and tree
size structure, and composes a side-by-side comparison figure for site~978
(S.~Appalachia)---one of the paper's ten representative sites that falls inside the
Tennessee box. The reproduced panel matches the paper's published panel~\#7 in both the
year-1000 diameter distribution (dominated by the $0$--$8$ and $8$--$28$~cm classes,
$\sim$45\% each, with a rapid tail-off) and the successional species composition
(a pitch pine, \textit{PINUrigi}, base with tulip poplar, sugar maple, white pine and
white ash accumulating over time). Because the seed is fixed and the RNG is
partition-invariant, the result is independent of the number of GPUs, substantiating
the determinism claim in $C_3$.

\arttime

\begin{description}
\item[Artifact Setup] $\approx$4~min: data load and directed-graph/METIS partitioning
  ($<$1~s), host-side agent build (\texttt{build\_from\_local\_data}, $\sim$192~s), and
  one-time GPU kernel generation ($\sim$32~s).
\item[Artifact Execution] $\approx$14.5~min (869.8~s) for the 1000-year simulation on
  2$\times$ P100. The first reported 50-year interval includes a one-time CUDA
  JIT/kernel-compile warmup ($\sim$330~s); steady-state 50-year intervals take
  $\sim$27~s each (compute plus snapshot I/O).
\item[Artifact Analysis] $\approx$3~min: extract site~978's full time series across the
  20 snapshots and render/compose the comparison figure.
\end{description}
Total wall-clock $\approx$22~min on the reference two-GPU node.

\artin

\artinpart{Hardware}

The reference reproduction node has 2$\times$ NVIDIA Tesla P100-SXM2 GPUs (16~GB each),
56 CPU cores and 125~GB host RAM (driver 560.35.05, CUDA 12.6). Requirements: at least
one CUDA (or ROCm) GPU with $\geq$16~GB; \textbf{two GPUs are recommended}---the 54
sites split $\sim$27 per GPU ($\sim$15M tree-agents each), which fits a 16~GB card,
whereas all 54 sites on a single 16~GB GPU ($\sim$30M agents) may exhaust memory. No
special interconnect is required (single node). The paper's production runs use OLCF
Frontier (AMD MI250X, ROCm); the same artifact runs there via the ROCm CuPy build.

\artinpart{Software}

\begin{itemize}
\item GGap --- this repository (branch \texttt{master}). URL: \texttt{https://github.com/TODO/GGap}.
\item SAGESim --- GPU agent-based framework, commit $\geq$\texttt{16cb647} (branch
  \texttt{dev}). URL: \texttt{https://github.com/TODO/SAGESim}. Installed via
  \texttt{pip install -e .}; provides \texttt{numpy}, \texttt{networkx}, \texttt{awkward}.
\item Python $\geq$3.13 (validated on 3.14.6).
\item CuPy --- hardware-matched: \texttt{cupy-cuda12x}~14.1.1 (CUDA~12.x) or a ROCm build.
\item \texttt{mpi4py}~4.1.2 with an MPI runtime (validated: OpenMPI~5.0.10).
\item \texttt{numpy}~2.5.1, \texttt{networkx}~3.6.1, \texttt{awkward}~2.11.0.
\item METIS~5.1.0 (\texttt{gpmetis} binary) --- required only for multi-GPU partitioning.
\item Figure step only: \texttt{matplotlib}, \texttt{pandas}, \texttt{pillow}.
\end{itemize}
\texttt{cupy} and \texttt{mpi4py} are intentionally not pinned by SAGESim because they
depend on the local CUDA/ROCm and MPI stack; reviewers install versions matching their
hardware.

\artinpart{Datasets / Inputs}

All inputs ship with the repository at
\texttt{SC2026/conus\_simulation/input\_data/} (seven \texttt{CONUS\_*.csv} files:
site, rangelist, climate, climate\_stddev, altitudes, specieslist, species\_sources); no
external download is needed. The artifact selects the Tennessee sub-region by bounding
box (latitude $34.9$--$36.7$, longitude $-90.5$ to $-81.5$), yielding 54 sites and 235
species. SHA-256 checksums of the input files are recorded with the artifact.

\artinpart{Installation and Deployment}

Install SAGESim (\texttt{pip install -e .}), then the hardware-matched
\texttt{cupy}/\texttt{mpi4py}, and (for a multi-GPU run) the \texttt{gpmetis} binary
(\texttt{conda install -c conda-forge metis}). Clone GGap; no \texttt{pip install} of
GGap itself is required---\texttt{run.py} puts the \texttt{gap} package on the path
automatically. If the local MPI is not CUDA-aware, leave
\texttt{MPICH\_GPU\_SUPPORT\_ENABLED} unset. Full step-by-step instructions are in
\texttt{SC2026/tn\_example/README.md}.

\artcomp

The workflow is a four-task pipeline, $T_1 \rightarrow T_2 \rightarrow T_3 \rightarrow T_4$:
\begin{description}
\item[$T_1$ Partition] Load the CONUS CSVs, apply the Tennessee bounding box (54 sites),
  build the directed seed-dispersal graph, and METIS-partition the sites across the MPI
  ranks (balanced 27/27 for two ranks). Number of partitions $=$ number of MPI ranks
  $=$ number of GPUs.
\item[$T_2$ Simulate] \texttt{mpirun -n 2 python run.py} --- run 1000 simulation years,
  writing a snapshot every 50 years plus per-rank metadata. Each rank binds to its own
  GPU. Output: \texttt{results/snapshots/}.
\item[$T_3$ Extract] \texttt{python extract\_timeseries.py} --- read all 20 snapshots for
  site~978 and write a per-year \texttt{species\_data.csv} time series.
\item[$T_4$ Figure] \texttt{python plot\_tn\_site978.py} then
  \texttt{python compose\_comparison.py} --- render the two panels and compose the
  side-by-side comparison against the paper's published figure.
\end{description}
Experimental parameters (defaults, reproducing the reference run): dispersal factor
$2.0$, 500 gaps/site, 1000 tree slots/gap, 1000 years, snapshot interval 50 years,
random seed 42 (identical on all ranks). The run is deterministic, so a single
replicate suffices; the seed makes 1-GPU and 2-GPU runs bitwise identical.

\artout

$T_2$ writes raw \texttt{.npz} snapshots and per-rank metadata; $T_3$ writes
\texttt{results/species\_timeseries/site\_0978/species\_data.csv} (per species, per
report year: six diameter-class counts, basal area, biomass~C/N); $T_4$ writes
\texttt{figures/tn\_site978\_vs\_paper.\{png,pdf\}}. The figure places the paper's
panel~\#7 above the reproduced panels; agreement in the year-1000 size distribution and
in the dominant successional species constitutes a successful reproduction.


% %%%%%%%%%%%%%%%%%%%%%%%%%%%%%%%%%%%%%%%%%%%%%%%%%%%%%
% %  AE Appendix
% %%%%%%%%%%%%%%%%%%%%%%%%%%%%%%%%%%%%%%%%%%%%%%%%%%%%%
\newpage
\appendixAE

\arteval{1}
\artin

Provision a node with $\geq$2 CUDA GPUs ($\geq$16~GB each; the reference used 2$\times$
Tesla P100). Then:
\begin{enumerate}
\item Install SAGESim: \texttt{git clone <SAGESim>; cd SAGESim; pip install -e .}
\item Install hardware-matched runtime packages:
  \texttt{pip install cupy-cuda12x mpi4py} (swap the CuPy build to match your CUDA/ROCm;
  \texttt{mpi4py} needs an MPI runtime such as OpenMPI installed first).
\item Install the METIS partitioner (multi-GPU only):
  \texttt{conda install -c conda-forge metis}.
\item Install figure dependencies: \texttt{pip install matplotlib pandas pillow}.
\item Clone GGap and enter the artifact directory:
  \texttt{git clone <GGap>; cd GGap/SC2026/tn\_example}.
\end{enumerate}

\artcomp

From \texttt{GGap/SC2026/tn\_example}, run the pipeline:
\begin{verbatim}
# T2: simulate (2 GPUs; use `python run.py`
#     for a single GPU)
mpirun -n 2 python run.py

# T3: extract site-978 full time series
python extract_timeseries.py

# T4: plot and compose the comparison figure
python plot_tn_site978.py
python compose_comparison.py
\end{verbatim}
Each step depends on the previous one. \texttt{run.py} accepts overrides
(e.g.\ \texttt{--years}, \texttt{--report\_interval}); the defaults reproduce the
reference run. (\texttt{python extract.py} is an optional convenience that writes
final-year tables for all 54 sites and is not part of the figure pipeline.)

\artout

The pipeline produces \texttt{figures/tn\_site978\_vs\_paper.png}. To evaluate:
compare the reproduced panels (bottom) with the paper's published panel~\#7 (top). The
reproduction is successful when (i) the year-1000 diameter distribution is concentrated
in the two smallest size classes ($0$--$8$ and $8$--$28$~cm) with a rapid tail-off, and
(ii) the species-composition-over-time panel shows the same dominant assemblage
(\textit{PINUrigi} base, with \textit{LIRItuli}, \textit{ACERsacc}, \textit{PINUstro},
\textit{FRAXamer}). This corresponds directly to the paper's Figure element for the ten
representative sites and substantiates the ecological-validity and determinism claims
($C_3$), while exercising the GPU model ($C_1$) and its multi-GPU distribution ($C_2$).
Note: the reference artifact saves snapshots every 50 years (20 points); the paper's
production figure uses a 10-year interval (100 points), so the reproduced composition
curve is coarser but follows the same trajectory.

\end{document}
